\section{Introduction}
\label{sec:introduction}
\vspace{-0.1cm}
Recent advances in 3D reconstruction~\citep{park2019deepsdf, sun2021neuralrecon, sun2022neuconw} and neural rendering~\citep{mildenhall2020nerf, sun2022direct, zhang2022nerfusion} have significantly enhanced the quality of 3D digitization of large scenes, such as buildings and city landmarks~\citep{hardouin2020next, zhang2021continuous, liu2022learning, xiangli2022bungeenerf, li2023matrixcity}.
However, image-capturing process still remains time-consuming and labor-intensive. To scan a building-scale scene using a drone, it may take several days' effort of a professional team to ensure full coverage. Moreover, even professional pilots may miss some areas at the first scan, leading to multiple rounds for rescanning.


To alleviate the effort in manual scanning, active 3D reconstruction algorithms, especially the Next-Best-View (NBV) policy, have emerged as a promising approach to automate view planning. 
However, typical rule-based NBV policies~\citep{lee2022uncertainty, pan2022activenerf, zhan2022activermap, peralta2020next, guedonscone, ran2023neurar} heavily rely on empirically designed view-selection criteria and hand-crafted action space, such as a hemisphere as shown in Fig.~\ref{fig:teaser}, which limit their generalizability to unseen scenes.
To enhance generalization capacity, pioneering work on learning-based NBV policy has been proposed~\citep{peralta2020next, zhan2022activermap, chenlearning, chaplotlearning}. Nevertheless, they are constrained into impractical setups in real-world scenarios, such as object-centric capturing and limited action space for ground robots, resulting in incomplete 3D reconstruction because of significant self-occlusion.


Therefore, in this work, we study this problem: \textit{``How to design an NBV policy that supports exploration in any space and also can generalize to unseen building-scale geometries?"} 
Building upon the design of previous reinforcement learning-based (RL-based) NBV policies~\citep{peralta2020next, chenlearning, chaplotlearning}, our setup faces the following new challenges that were not fully studied before:
1) Designing an easy-to-generalize action space, rather than object-specific spaces used in previous work, such as hand-crafted candidates or a constrained hemisphere;
2) With a large action space, an informative representation is necessary to guide the policy to efficiently find an optimal capturing trajectory;
3) Proposing generic reward functions and developing the distributed training procedure for better generalizability.

To overcome these two challenges, we propose \textit{GenNBV}, an end-to-end generalizable NBV policy. It has a novel design of action spaces and state embeddings, which allows free-space exploration and ensures a large coverage when applied to unseen objects.
We first extend the former limited action space, like a hemisphere, to 5D free space. As shown in Fig.~\ref{fig:teaser}, this new space is composed of a position subspace of approximately $20m$ x $20m$ x $10m$ and an omnidirectional heading subspace. This free space enables our agent drone to scan from any viewpoint.
Moreover, with a much larger action space, our RL-based policy can be easily generalized to building-scale geometries without concerning the potential scale shift from training to evaluation. In contrast, many hand-crafted designs~\citep{pan2022activenerf, lee2022uncertainty} limit the available viewpoints to a small predefined set, and thus cannot generalize if there is a drastic change in scales.


Second, in order to generalize a well-trained NBV policy to unseen scenes during evaluation, we propose a novel state representation directly from raw sensory inputs to guide the policy. In contrast to the widely used neural radiance field (NeRF)~\citep{pan2022activenerf, zhan2022activermap, lee2022uncertainty, smith2022uncertainty,ran2023neurar}, our representation does not need time-consuming per-scene optimization. 
Specifically, our NBV policy incorporates a multi-source state embedding, which comprises three key components: a novel geometric embedding, a semantic embedding, and an action embedding. The geometric embedding is derived from multi-view depth maps and represents an encoded probabilistic 3D grid, while the semantic embedding is extracted from RGB images, and the action embedding is an encoded viewpoint sequence.
In contrast to 2D representations used in previous RL-based NBV policies~\citep{chenlearning, chaplotlearning}, our multi-source state embedding, extracted from a probabilistic 3D grid, is a more comprehensive representation that supports a robust prediction of next viewpoints.


% Our Experiments
To validate the effectiveness of GenNBV, we construct a benchmark for training and evaluation with Houses3K~\citep{houses3k} dataset from the NVIDIA Isaac Gym~\citep{makoviychuk2021isaac} simulator. We further test its generalization ability on novel buildings and object categories from OmniObject3D~\citep{wu2023omniobject3d} and Objaverse~\citep{deitke2023objaverse} datasets.
Our evaluation metrics quantify the completeness, accuracy, and efficiency of active 3D reconstruction.
To consolidate completeness and efficiency into a single number, we also propose a metric, the Area Under the Curve (AUC) of the coverage. 
Compared to heuristic~\citep{isler2016information}, information gain-based~\citep{isler2016information, lee2022uncertainty, zhan2022activermap}, and RL-based baselines~\citep{peralta2020next}, our method achieves the best result under different metrics.
We also provide visualizations to demonstrate the generalizability of GenNBV.
A demo video is attached in the appendix to further illustrate our method.
